\section{The mixed transition in \texorpdfstring{$X_9$}{X9}}
\label{sec:examples}\label{subsec:x9}

\subsection{Geometry and the smoothing}

Let $S_7=\operatorname{Bl}_{p_1,p_2}\mathbb P^2$ denote the del Pezzo
surface of degree seven. The compact threefold $X_9$ has two ordinary
double points and one cone over $S_7$. Its crepant resolution
$\wideX_9$ is an anticanonical hypersurface in the toric fourfold with
fan polytope
\begin{equation}
\begin{split}
\Delta_9=\operatorname{conv}\{&v_0=(1,0,0,0),\ v_1=(0,1,0,0),\ v_2=(-1,-1,0,0),\\
&v_3=(0,0,1,0),\ v_4=(0,0,0,1),\ v_5=(-4,-2,-1,0),\\
&v_6=(-4,-2,0,-1),\ v_7=(3,1,1,1),\ v_8=(2,1,1,1)\}.
\end{split}
\label{eq:x9-polytope}
\end{equation}
The singular two-faces are the squares $2367$, $2457$ and the
pentagon $34568$. We use the square diagonals $67$ and $24$, and
the star subdivision of the pentagon at $p=(-2,-1,0,0)$.
The exceptional divisor $E=D_p$ is $S_7$; the two exceptional
curves are $C_1,C_2$. The remaining boundary point $(-1,0,0,0)$
lies in a facet interior and its divisor misses the hypersurface.
The chamber and integral intersection data are specified in
Appendices~\ref{app:integral-charges} and~\ref{app:mirror-periods}.

Write $h,e_1,e_2$ for the standard basis of $H_2(E;\ZZ)$,
with intersection form $\operatorname{diag}(1,-1,-1)$ and
$K_E=-3h+e_1+e_2$. The compact root relation is
\begin{equation}
 R:=\iota_*(e_1-e_2)=C_1+C_2.
 \label{eq:x9-root-relation}
\end{equation}
Neither the nodal curves nor the root can deform independently.
In reduced local coordinates $(u_1,u_2,\beta)$ the allowed
direction is $(-1,-1,1)$, as the charge calculation below shows.

This direction integrates to an actual smoothing
$\mathcal X\to\Delta$ with smooth general fiber $X_{9,t}$
\cite[Theorem~B]{Wuebben:2026smoothing}. The construction uses the
facet $F=\{v_2,\ldots,v_8\}$ and deformation degree $(0,1,-1,-1)$.
On its ordered edges
\begin{equation}
 (25,26,27,36,37,38,45,47,48,56,78)
 \label{eq:x9-facet-edges}
\end{equation}
the dilation vector is $(0,0,1,1,0,0,1,0,0,0,0)$. It defines
compatible Minkowski decompositions into two segments on each square
and a segment plus a unimodular triangle on the pentagon. The
Ilten--Vollmert deformation and its relative anticanonical family
therefore smooth all three germs
\cite{Altmann:1995,Ilten:2012,Wuebben:2026smoothing}, with
\begin{equation}
 (u_1,u_2,\beta)=t(-1,-1,1).
 \label{eq:x9-family-direction}
\end{equation}
Smoothability here is an existence result, not an inference from
the first-order kernel.

The massless multiplets on the two smooth sides are
\begin{equation}
\begin{array}{c|rrr|rr}
 &h^{1,1}&h^{2,1}&\chi&n_V&n_H^{\rm neutral}\\\hline
 \wideX_9&6&122&-232&5&123\\
 X_{9,t}&3&123&-240&2&124 .
\end{array}
\label{eq:x9-spectrum-table}
\end{equation}
Here $n_V=h^{1,1}-1$ excludes the graviphoton and
$n_H^{\rm neutral}=h^{2,1}+1$ includes the universal hypermultiplet.
The general transition formula of Appendix~\ref{app:deformations}
gives the same change as the independent construction of the general
fiber: two flavor vectors are Higgsed, the rank-one $E_2$ Coulomb
vector disappears, and one new neutral hypermultiplet remains.

\subsection{Integral charges and the joint constraint}

The six divisors
\begin{equation}
 \mathcal D=(D_2,D_5,D_6,D_7,D_8,E)
 \label{eq:x9-divisor-basis}
\end{equation}
form the full free integral $H^2$ lattice, not merely a saturated
toric quotient. An explicit unimodular pairing proves this in
Appendix~\ref{app:integral-charges}. Their pairings with
$(C_1,C_2,R)$ are
\begin{equation}
 Q_{\mathcal D}=
 \begin{pmatrix}1&-1&0\\0&1&1\\-1&0&-1\\-1&1&0\\0&0&0\\0&0&0\end{pmatrix}.
 \label{eq:x9-full-matrix}
\end{equation}
An integral change of vector basis isolates
\begin{equation}
 V_1=-D_6,\qquad V_2=D_5,\qquad
 Q_{X_9}^{\rm eff}=\begin{pmatrix}1&0&1\\0&1&1\end{pmatrix}.
 \label{eq:x9-effective-vectors}
\end{equation}
Both nonzero Smith invariants are one. Thus
\begin{equation}
 u_1+\beta=u_2+\beta=0,\qquad
 \ker Q_{X_9}^{\rm eff}=\CC(-1,-1,1).
 \label{eq:x9-moment-maps}
\end{equation}
Every nonzero allowed deformation moves all three sectors. The
exceptional-divisor vector is neutral under this flavor constraint:
$E|_E=K_E$ and $K_E\cdot(e_1-e_2)=0$. Its loss in
\eqref{eq:x9-spectrum-table} is the loss of the local Coulomb
direction, not a third row of the rank-two flavor matrix.

The reduced rigid current ring gives a useful refinement of this
coordinate statement. If $J_\pm,J_0$ are the $E_2$ instanton currents
and $u_i=h_i\widetilde h_i$, its invariant generators can be chosen as
\begin{equation}
 X=h_1h_2J_-,\qquad Y=\widetilde h_1\widetilde h_2J_+,\qquad
 XY=\beta^4.
 \label{eq:x9-dressed-generators}
\end{equation}
Consequently the formal reduced quotient is $\CC^2/\ZZ_4$.
The integral extension of the flavor action and the current-ring
calculation are given in Appendix~\ref{app:integral-charges}.
This is neither a finite-Planck hypermultiplet metric nor an
unbroken discrete gauge group at a generic Higgs vacuum.

\subsection{Surviving vectors and classical couplings}

Let $Y$ be the common complement of small singular neighborhoods.
The nodal Milnor fibers have $H^2=0$. For the $S_7$ cone the Milnor
fiber is $\operatorname{Bl}_{pt}\mathbb P^3\setminus S_7$;
its restriction to the link has image $R^\perp/\ZZ K_E$.
Integral Mayer--Vietoris gluing therefore gives
\begin{equation}
 H^2(X_{9,t};\ZZ)_{\rm free}
 \simeq \ker(C_1,C_2,R)/\ZZ E.
 \label{eq:x9-smooth-h2}
\end{equation}
A primitive basis of this quotient is
\begin{equation}
 (A,B,\nu)=(D_2+D_7,D_5+D_6-D_7,\nu),\qquad
 D_0=A+4B-2\nu,\quad D_1=A+2B-\nu,
 \label{eq:x9-primitive-smooth-basis}
\end{equation}
where $\nu$ agrees with $D_8$ on $Y$. The older convenient basis
$(D_0,D_1,\nu)$ has index two and will be used only when that
particular comparison simplifies the geometry.

The surface $D_8\cong\mathbb F_1$ meets $E$ in its exceptional
curve. In the smoothing it becomes a rigid plane of class $\nu$
and normal bundle $\mathcal O_{\mathbb P^2}(-3)$, hence a local
$E_0$ geometry \cite{Morrison:1996pp}. The cut-and-paste calculation
in Appendix~\ref{app:integral-charges} gives
\begin{equation}
 \nu^3=9\quad\hbox{versus}\quad D_8^3=8,\qquad
 c_2(X_{9,t})\cdot\nu=-6\quad\hbox{versus}\quad
 c_2(\wideX_9)\cdot D_8=-4.
 \label{eq:x9-smooth-cubic}
\end{equation}
All cubic entries involving $D_0$ or $D_1$ agree when these fixed
representatives on $Y$ are used. In particular,
\begin{equation}
 \mathcal F(\alpha D_0+\beta D_1)
 =\frac{83\alpha^3+147\alpha^2\beta+81\alpha\beta^2+14\beta^3}{6}
 \label{eq:x9-inherited-cubic}
\end{equation}
on both sides, with $\mathcal F=J^3/6$ in five-dimensional units.

For the quantum comparison a different lift is required: it must
annihilate all four electric charges, not just their root difference.
In the primitive nef basis $H_1,\ldots,H_6$ of
\eqref{eq:x9-mirror-nef}, this lattice is
\begin{equation}
 \Lambda^0_{\rm mag}
 =\{P:P C_1=P C_2=P\iota_*e_1=P\iota_*e_2=0\}
 =\ZZ H_2\oplus\ZZ H_3\oplus\ZZ H_4.
 \label{eq:x9-unconfined-magnetic-lattice}
\end{equation}
The classes in \eqref{eq:x9-primitive-smooth-basis} have lifts
\begin{equation}
 \widetilde A=D_2+D_7,\qquad
 \widetilde B=D_5+D_6-D_7+E,\qquad \widetilde\nu=D_8+E,
 \label{eq:x9-unconfined-lifts}
\end{equation}
whose columns in $(H_2,H_3,H_4)$ form
\begin{equation}
 U=\begin{pmatrix}0&-1&-2\\2&-1&-1\\-1&1&1\end{pmatrix},\qquad \det U=-1.
 \label{eq:x9-unconfined-unimodular}
\end{equation}
Indeed $P\mapsto P+(P\cdot\iota_*e_2)E$ induces the integral
isomorphism from \eqref{eq:x9-smooth-h2} to
\eqref{eq:x9-unconfined-magnetic-lattice}. With these lifts all ten
cubic entries and all three second-Chern pairings agree with the
smooth side. In lexicographic order in $(A,B,\nu)$ they are
\begin{equation}
 (-1,1,4,-1,-4,-6,3,4,6,9),\qquad (14,30,-6).
 \label{eq:x9-primitive-cubic}
\end{equation}
The equality for unconfined lifts and the jump in
\eqref{eq:x9-smooth-cubic} compare different resolved representatives.
The $E$ shifts reconcile them. Section~\ref{sec:quantum} determines
whether this integral polynomial comparison extends to quantum couplings.

The integral lattice also screens the local $E_0$ line-charge group.
A plane line has $(A,B,\nu)$ charges $(2,-2,-3)$, whereas the compact
class $D_1^2$ has charges $(1,7,1)$. Thus the local $\ZZ_3$ electric
one-form symmetry does not persist in the compact charge lattice
\cite[Sections~3.4 and~5.3]{Morrison:2020ool}. This uses integral
combinations of effective curves; it does not assert an irreducible
BPS representative of $D_1^2$.

\subsection{A wrapped-M5 anomaly check}
\label{subsec:x9-m5-anomaly}

Wrap an M5-brane on $S_-=D_8\cong\mathbb F_1$ and compare it with
the string on $S_+=\nu\cong\mathbb P^2$. This comparison fixes the
representative $D_8$, rather than its unconfined lift $D_8+E$.
The restrictions of $(D_0,D_1,D_8)$ to
$\mathbb F_1=\operatorname{Bl}_{pt}\mathbb P^2$ are $(0,h,-3h+e)$,
with $h^2=1,e^2=-1,h e=0$; on the plane they are $(0,h,-3h)$.
The M5 tensor lattice consequently loses one negative direction
\cite{Maldacena:1997de}.

In the inflow conventions of \cite[eq.~(2.21)]{Katz:2020ewz},
\begin{equation}
 I_4(S)=-\frac12(SD_AD_B)F^AF^B
 -\frac{S^3+\tfrac12c_2(X)S}{6}c_2(N)
 +\frac{c_2(X)S}{48}p_1(T\Sigma),
 \label{eq:x9-string-anomaly-general}
\end{equation}
where $N$ denotes the transverse $SU(2)$ bundle. The two transverse
coefficients are both $-1$, while
\begin{equation}
 I_4(S_+)-I_4(S_-)
 =-\frac12(F^\nu)^2-\frac1{24}p_1(T\Sigma).
 \label{eq:x9-string-anomaly-jump}
\end{equation}
The removed left-moving lattice boson has unit $\nu$ charge and
accounts for both terms. The exceptional direction approaches $e$
in the blowdown limit; it need not be anti-self-dual at every
polarization. Thus worldsheet content and bulk Chern--Simons inflow
change together. This protected anomaly check does not identify an
interacting infrared central charge or count the Higgs condensates.
